\section{The Cusp Is Physical Space}

The geometry part of this paper has built a four-dimensional crystal, the hyperbolic honeycomb $\{3,4,3,4\}$ whose cells are $24$-cells. That bulk is where the directions, the spin, and the gauge structure live. It is not, however, the space we walk around in. This section makes the central claim that fixes what space is. Physical three-dimensional space is the flat $\{4,3,4\}$ cubic \emph{cusp} of that bulk, the boundary layer that sits at an ideal corner of the crystal. Space is not a fifth ingredient added by hand. It is the radiance cast by the burning bulk, the flat three-dimensional skin of the four-dimensional weave.

\begin{principle}[Space is the boundary, time is the growth]
Physical space is the flat $\{4,3,4\}$ cubic cusp of the $\{3,4,3,4\}$ bulk, taken at one beat. Time is the growth of the mesh, the wake receding outward beat by beat. The $3+1$ of physics is three cusp directions plus the growth beat.
\end{principle}

Before formalizing, two terms need plain statements, since the whole section turns on them. A \emph{bulk} is the full four-dimensional interior of the honeycomb, the stack of $24$-cells meeting at their faces. A \emph{cusp} is the flat infinite tiling that lives at an ideal vertex of a paracompact honeycomb, the surface a horosphere meets as it touches the boundary at infinity. The reader who has met the bulk and the staircase in the two preceding sections can take the cusp as the flat layer those parabolic reflections build.

\subsection{The claim in one line}

The model has three layers, and they answer three different questions. The bulk answers where the internal structure lives. The cusp answers where matter moves. The growth answers when.

\begin{table}[ht]
\centering
\begin{tabular}{@{}ll@{}}
\toprule
layer & what it is \\
\midrule
the bulk & the 4D hyperbolic honeycomb of 24-cells, where spin and gauge live \\
the cusp & the flat $\{4,3,4\}$ cubic horosphere at an ideal corner, our 3D space \\
time & the growth of the mesh, the wake receding outward \\
\bottomrule
\end{tabular}
\caption{The three layers of the model and the role each one plays.}
\label{tab:three-layers}
\end{table}

Space is not a separate fourth ingredient grafted onto the eight atoms. It is the \emph{boundary} of the bulk crystal, the flat three-dimensional cubic layer at an ideal corner. The mesh atom is the whole $\{3,4,3,4\}$ weave. What we experience as space is one particular flat slice of that weave, the cubic cusp. The honeycomb section and the staircase section built exactly this object, the parabolic reflections that close the corner into a flat $\{4,3,4\}$ tiling.

\subsection{Paracompactness is a feature}

A reader who knows the classification of hyperbolic honeycombs may worry about a word attached to $\{3,4,3,4\}$. It is \emph{paracompact}. Its vertex figure is the Euclidean $\{4,3,4\}$ cubic honeycomb, which is an infinite ideal element rather than a finite one. In the older selection language this sounded like a flaw to be apologized for. It is the opposite.

That infinite ideal element \emph{is} the flat three-dimensional physical space. The paracompactness is not a defect to be explained away. It is exactly what supplies a clean, periodic, flat three-dimensional crystal for physics to live on. A genuinely compact honeycomb, one with finite cells all the way down, would have no such flat cubic layer to offer.

\begin{theorem}[The cusp is the cubic honeycomb]
The vertex figure of the paracompact honeycomb $\{3,4,3,4\}$ is the Euclidean cubic honeycomb $\{4,3,4\}$. The ideal vertex therefore carries an exactly flat, periodic, three-dimensional cubic tiling. This tiling is the physical space of the model.
\end{theorem}

There is a trade hidden in this choice, and it is worth stating squarely. The genuinely compact honeycombs, the ones with finite cells, are all \emph{non-crystallographic}. They lack the spinor on their sites, the half-integer-spin pieces that the matter sector needs. So the two requirements pull against each other. Either you accept ideal vertices and get a clean cubic cusp with a spinor on the sites, or you take a compact honeycomb and lose the spinor. The substrate takes the cusp.

\begin{result}[Paracompactness is the price and the prize]
Paracompactness is what gives the $D_4$ cusp its flat cubic space and its spinor sites at once. The compact alternatives are non-crystallographic and carry no spinor. The substrate accepts ideal vertices on purpose, because that is what buys both the flat space and the spin.
\end{result}

\subsection{Flat at finite distance, not only at infinity}

A natural worry follows. If the cusp lives at an ideal vertex, at the boundary at infinity, is the flat space real only in the infinite limit? Would a finite observer ever stand on a truly Euclidean surface, or only approach one?

The answer is that the flat space is real at every finite level. A horosphere in hyperbolic space is intrinsically \emph{Euclidean at every finite distance}, not merely in a limit. This is a classical theorem about horospheres, and the model confirms it by direct measurement. The signature is how much room a band of the surface gains as you step outward.

\begin{table}[ht]
\centering
\begin{tabular}{@{}ll@{}}
\toprule
surface & growth of a band \\
\midrule
the bulk & exponential ($\sim 9$) \\
a horosphere & polynomial ($\sim 1.5$) \\
\bottomrule
\end{tabular}
\caption{Band growth on the bulk versus on a horosphere. The bulk gains room exponentially, the horosphere only polynomially, the signature of an intrinsically flat surface.}
\label{tab:band-growth}
\end{table}

A band of the bulk gains room exponentially, growing by roughly a factor of nine per step, the hallmark of negative curvature. A band of the horosphere gains room polynomially, by roughly a factor of one and a half, exactly as a flat Euclidean shell does. The flatness is present at finite radius, not deferred to infinity.

\begin{result}[The flat space is finite and exactly Euclidean]
A horosphere band of the $\{3,4,3,4\}$ bulk grows polynomially, not exponentially. The flat three-dimensional space is a real, finite-distance, exactly Euclidean surface, not an object that exists only in a limit. This is a measured emergent result at depth level L1 on the cubic-cusp geometry.
\end{result}

\subsection{The clean cusp versus a generic slice}

One more guard is needed. Not every slice of the bulk is physical space. A reader might suppose that any horosphere, sliced anywhere, would do. It would not. Only the clean cubic cusp is space, and the difference between it and a generic slice is measurable.

The right measure is the \emph{spectral dimension}, the effective dimensionality a random walk feels as it explores the surface. A clean three-dimensional crystal reads close to three. A rough, aperiodic surface reads lower, because the walk keeps getting trapped and the surface behaves as if it had fewer than three directions.

\begin{table}[ht]
\centering
\begin{tabular}{@{}lll@{}}
\toprule
slice & spectral dimension & character \\
\midrule
clean $\{4,3,4\}$ cubic cusp & $2.97$ & clean 3D \\
generic horosphere slice & $2.16$ & aperiodic, rough, sub-3D \\
\bottomrule
\end{tabular}
\caption{Spectral dimension of the clean cubic cusp versus a generic horosphere slice. Only the clean cusp reads as genuine three-dimensional space.}
\label{tab:spectral-dimension}
\end{table}

The clean cubic cusp reads a spectral dimension of $2.97$, essentially three, a clean three-dimensional space. A generic horosphere slice is aperiodic and rough, reading only $2.16$, distinctly sub-three-dimensional. So the clean cubic cusp, and not a disordered generic slice, is physical space.

\begin{result}[Only the clean cubic cusp is 3D space]
The clean $\{4,3,4\}$ cubic cusp has spectral dimension $2.97$, essentially three. A generic horosphere slice is aperiodic and rough, with spectral dimension $2.16$, sub-three-dimensional. Physical space is the clean cubic cusp, not an arbitrary slice. This is a measured result on the cusp geometry.
\end{result}

An earlier and more optimistic reading had hoped that an aperiodic slice might still be isotropic, that disorder might average into rotational symmetry for free. That reading did not survive the test. The clean cubic cusp, with its periodic cubic order, is the right physical space, and the rough slice is not.

\subsection{The bulk-cusp interface is a projection}

Spin and gauge live in the four-dimensional bulk sites, the $24$ directions of the $24$-cell, the $D_4$ roots. Matter, on the other hand, propagates on the three-dimensional cubic cusp. How are the two related? Not by an unbridged gap that must be papered over. By a \emph{projection}, the standard dimensional reduction of a slice.

\begin{principle}[The interface is a projection, not a gap]
The bulk and the cusp are joined by projection along the radial direction that the horosphere picks out. Bulk structures, the spin and the directions and the internal $so(8)$, project onto the cusp tangent. Nothing is severed at the interface, the bulk data simply casts a three-dimensional shadow.
\end{principle}

The rest of this section follows three of those bulk structures across the interface and reads off their shadows. The spinor branches, the directions reduce, and the internal symmetry rides along as quantum numbers.

\subsection{The spinor branch, SO(4) to SO(3)}

The bulk's spacetime rotation group is $\mathrm{SO}(4)$, which factors as $\mathrm{SU}(2) \times \mathrm{SU}(2)$. This is the rotation symmetry of the four-dimensional crystal before any slice is chosen.

Choosing a horosphere picks out a radial direction, the direction along which the cusp sits at infinity. That choice \emph{breaks} $\mathrm{SO}(4)$ down to $\mathrm{SO}(3)$, the three-dimensional rotation group of the cusp, realized as the diagonal $\mathrm{SU}(2)$ inside the product. Under this branch the four-dimensional Dirac spinor $(2,1) + (1,2)$ decomposes to $2 + 2$, two three-dimensional Pauli spinors, each carrying spin one half.

\begin{table}[ht]
\centering
\begin{tabular}{@{}lll@{}}
\toprule
step & symmetry & spinor \\
\midrule
bulk & $\mathrm{SO}(4) = \mathrm{SU}(2) \times \mathrm{SU}(2)$ & 4D Dirac, $(2,1) + (1,2)$ \\
pick a horosphere & break to $\mathrm{SO}(3)$, diagonal $\mathrm{SU}(2)$ & $4$ branches to $2 + 2$ \\
cusp & $\mathrm{SO}(3) = \mathrm{SU}(2)$ & two 3D Pauli spinors, spin $1/2$ each \\
\bottomrule
\end{tabular}
\caption{The spinor branching across the bulk-cusp interface. The four-dimensional Dirac spinor reduces to two three-dimensional Pauli spinors.}
\label{tab:spinor-branch}
\end{table}

This branching was verified directly. The quadratic Casimir came out as $3/4$ on each piece, the exact value $s(s+1)$ for $s = 1/2$, confirming that each reduced spinor genuinely carries spin one half. So the four-dimensional bulk Dirac spinor \emph{projects} to three-dimensional Pauli spinors. The bulk spin becomes the cusp spin.

\begin{proposition}[Cusp spin is the shadow of bulk spin]
Under the breaking $\mathrm{SO}(4) \to \mathrm{SO}(3)$ induced by choosing a horosphere, the 4D Dirac spinor $(2,1)+(1,2)$ branches to $2+2$, two 3D Pauli spinors. The Casimir is $3/4$ on each, so each carries spin $1/2$. This is the textbook reduction of a $(3+1)$D Dirac fermion to its 3D spatial spinor, $\mathrm{Spin}(3) = \mathrm{SU}(2)$. It is a derived structural result at depth level L2.
\end{proposition}

This dissolves what could have looked like a mismatch. The spinor lives on the bulk sites, the matter we observe lives in three-dimensional space, and one might fear the two simply do not fit. They fit exactly. The spinor sites and physical space are not at odds. The cusp spin is the shadow of the bulk spin, the same object seen one dimension down.

\subsection{The directions project}

The $24$ bulk directions also cast a shadow on the cusp. Projecting them onto the cusp tangent space, orthogonal to the chosen radial direction, collapses the $24$ into a smaller set of distinct three-dimensional directions.

\begin{table}[ht]
\centering
\begin{tabular}{@{}ll@{}}
\toprule
set & count \\
\midrule
bulk sites & 24 directions \\
projected cusp sites & $\sim 9$ to $12$ distinct directions \\
naive cubic & 6 directions \\
\bottomrule
\end{tabular}
\caption{The $24$ bulk directions project onto the cusp tangent into roughly $9$ to $12$ distinct directions, richer than the naive cubic six.}
\label{tab:direction-projection}
\end{table}

The projected cusp sites give roughly nine to twelve distinct directions, \emph{richer} than the naive cubic six of a plain cubic lattice. They coarse-grain to the $\{4,3,4\}$ cubic continuum in the infrared. The practical consequence is welcome. Richer projected sites mean the cusp's residual anisotropy is even smaller than a naive cubic estimate would suggest, since more directions average more nearly into rotational symmetry.

\subsection{The symmetry chain and the anisotropy}

Pulling the threads together, the symmetry reduces in a clean chain along the interface. The bulk carries $so(8)$ on its sites and $\mathrm{SO}(4)$ as its spacetime rotation. Choosing the cusp drops the rotation symmetry to the cusp's octahedral cubic point group, a finite subgroup of $\mathrm{SO}(3)$. Full continuous $\mathrm{SO}(3)$ then re-emerges only in the infrared, after coarse-graining.

\begin{table}[ht]
\centering
\begin{tabular}{@{}lll@{}}
\toprule
layer & exact symmetry & isotropy \\
\midrule
bulk & $D_4$ / $F_4$ & isotropic to order 4 \\
cusp & finite octahedral & small order-4 anisotropy \\
infrared limit & $\mathrm{SO}(3)$ restored & isotropic \\
\bottomrule
\end{tabular}
\caption{The symmetry chain from bulk to cusp to infrared. The exact cusp symmetry is the finite octahedral group, and continuous rotation is emergent.}
\label{tab:symmetry-chain}
\end{table}

The cusp's exact symmetry is the finite octahedral group. This is why the cusp carries a small \emph{order-four anisotropy}, while the bulk, with its $D_4$ and $F_4$ symmetry, is isotropic all the way to order four. Full $\mathrm{SO}(3)$ is restored only in the infrared. The base symmetry is finite and discrete throughout. Continuous rotation is an emergent property of the long-wavelength limit, not a feature of the base.

\subsection{What rides along}

The spin reduces and the directions reduce, but the internal structure carries over intact. The internal $so(8)$ structure, which encodes the gauge content and the three generations through triality, rides along the interface as internal quantum numbers.

\begin{table}[ht]
\centering
\begin{tabular}{@{}ll@{}}
\toprule
bulk feature & cusp role \\
\midrule
4D Dirac spin & 3D Pauli spin, the cusp spin \\
24 directions & the projected cusp sites \\
$so(8)$ internal & gauge and generation quantum numbers \\
\bottomrule
\end{tabular}
\caption{What each bulk feature becomes on the cusp. The spin and directions reduce, and the internal symmetry carries over as quantum numbers.}
\label{tab:rides-along}
\end{table}

So the spinor reduces from four-dimensional Dirac to three-dimensional Pauli, the directions reduce from twenty-four to the projected cusp set, full $\mathrm{SO}(3)$ is restored in the infrared, and the $so(8)$ structure carries over as the internal quantum numbers that label gauge charge and generation.

\subsection{Time is the growth}

Space is the cusp taken at one beat. Time is something else entirely, and the model keeps the two apart on purpose. Time is the growth of the mesh.

The wake atom is the growing edge of the honeycomb, the frontier where new docks are born at peace each beat. The wake of the finite, ever-growing honeycomb recedes outward beat by beat. That receding edge is the \emph{present}. The growth is treated fully where finiteness and growth are taken up.

\begin{table}[ht]
\centering
\begin{tabular}{@{}ll@{}}
\toprule
dimension & what it is in the model \\
\midrule
the three spatial directions & the flat cubic cusp at one beat \\
time & the beat index, the growth of the mesh \\
\bottomrule
\end{tabular}
\caption{Space and time in the model. Space is the cusp geometry at one beat, time is the growth dynamics.}
\label{tab:space-and-time}
\end{table}

So the $3+1$ of physics is three cusp directions plus the growth beat. Space is geometry, time is dynamics. This is why space and time are separate atoms at the base. They are different kinds of thing, one a slice and one a count of slices, and their eventual unification into a single spacetime is itself an emergent fact taken up later.

\begin{principle}[Space is geometry, time is dynamics]
The three spatial directions are the flat cubic cusp at one beat. Time is the beat index, the growth of the mesh. The unification of the two into Lorentzian spacetime is emergent, which is exactly why the mesh and the beat are distinct atoms at the base.
\end{principle}

\subsection{Why it matters and the residual}

This section pins down what space \emph{is}. It is the clean $\{4,3,4\}$ cubic cusp, genuinely and cleanly three-dimensional, exactly flat at finite distance, and joined to the bulk by a clean projection rather than an unbridged gap. The three load-bearing results stack as follows.

\begin{result}[The cusp is physical 3D space]
Physical space is the clean $\{4,3,4\}$ cubic cusp of the $\{3,4,3,4\}$ bulk. The spin reduces $\mathrm{SO}(4)$ to $\mathrm{SO}(3)$, the 4D Dirac spinor branching to $2+2$ (derived, verified by the $3/4$ Casimir). The directions reduce to the projected cusp sites (measured). Full $\mathrm{SO}(3)$ is restored in the infrared. The space is exactly Euclidean at finite distance, with spectral dimension $2.97$.
\end{result}

One residual remains, and honesty requires naming it as permanent. The cusp carries a small cubic anisotropy inherited from its finite octahedral symmetry. This is a tiny imperfection that does not vanish with growth. It is a real, if small, prediction of the model rather than an artifact to be tuned away.

The remaining work past that residual is bookkeeping. The exact lattice-level cusp sites need to be enumerated, and the $so(8)$ carry-over into the cusp's internal quantum numbers needs to be written out in full. The mechanism for both is standard dimensional reduction, so the path is clear even where the details are not yet filled in.
